\documentclass[aspectratio=169]{beamer}
\usetheme{Madrid}
\usecolortheme{whale}

\usepackage[utf8]{inputenc}
\usepackage[spanish]{babel}
\usepackage{amsmath,amssymb}
\usepackage{graphicx}
\usepackage{listings}
\usepackage{booktabs}

\lstset{
    language=Python,
    basicstyle=\ttfamily\footnotesize,
    keywordstyle=\color{blue}\bfseries,
    stringstyle=\color{red},
    commentstyle=\color{gray}\itshape,
    breaklines=true,
    showstringspaces=false
}

\title[INT210 - Sesión 0.1]{Sesión 0.1: Bifurcaciones Lógicas y Decisiones Condicionales}
\subtitle{Módulo 0: Fundamentos Algorítmicos en Python}
\author[Diplomado en Ciencia de Datos]{Cátedra de Inteligencia Artificial y Ciencia de Datos}
\institute[INT210]{Machine Learning From Scratch}
\date{Semana 1}

\begin{document}

\frame{\titlepage}

\begin{frame}{Contenido de la Sesión}
    \tableofcontents
\end{frame}

\section{Introducción y Analogía}

\begin{frame}{La Analogía Infantil: El Guardián del Parque de Diversiones}
    \begin{columns}
        \column{0.55\textwidth}
        \begin{itemize}
            \item \textbf{Condición Principal (\texttt{if}):} Estatura $\ge 140$ cm y boleto dorado $\rightarrow$ Acceso a primera fila.
            \item \textbf{Alternativa (\texttt{elif}):} Estatura $\ge 120$ cm con tutor $\rightarrow$ Acceso supervisado.
            \item \textbf{Caso por defecto (\texttt{else}):} Redirección a juegos infantiles.
        \end{itemize}
        \column{0.45\textwidth}
        \begin{alertblock}{Principio Fundamental}
            El procesador evalúa proposiciones lógicas en secuencia determinista y bifurca el flujo de ejecución hacia una única rama válida.
        \end{alertblock}
    \end{columns}
\end{frame}

\section{Formalismo Matemático: Álgebra de Boole}

\begin{frame}{Álgebra Booleana y Tablas de Verdad}
    Para dos proposiciones lógicas $P, Q \in \{\text{True}, \text{False}\}$:
    \begin{center}
    \begin{tabular}{ccccc}
        \toprule
        $P$ & $Q$ & $P \land Q$ (\texttt{and}) & $P \lor Q$ (\texttt{or}) & $\neg P$ (\texttt{not}) \\
        \midrule
        True & True & \textbf{True} & \textbf{True} & \textbf{False} \\
        True & False & \textbf{False} & \textbf{True} & \textbf{False} \\
        False & True & \textbf{False} & \textbf{True} & \textbf{True} \\
        False & False & \textbf{False} & \textbf{False} & \textbf{True} \\
        \bottomrule
    \end{tabular}
    \end{center}
    \vspace{0.3cm}
    \textbf{Evaluación de Cortocircuito (\textit{Short-Circuit}):}
    \begin{itemize}
        \item En \texttt{A and B}: Si $A = \text{False}$, $B$ jamás es evaluado.
        \item En \texttt{A or B}: Si $A = \text{True}$, $B$ se omite por completo.
    \end{itemize}
\end{frame}

\section{Implementación y Casos de Estudio}

\begin{frame}[fragile]{Caso de Estudio 1: Firewall por Reglas en Ciberseguridad}
\begin{lstlisting}
def clasificador_firewall(ip_origen, puerto, protocolo, tamano_bytes):
    if ip_origen in LISTA_NEGRA_IPS:
        return 'Ataque'
    elif puerto in PUERTOS_CRITICOS:
        return 'Ataque'
    elif protocolo == 'UDP' and tamano_bytes > 1500:
        return 'Ataque'
    else:
        return 'Normal'
\end{lstlisting}
\begin{exampleblock}{Ventaja Arquitectónica}
    Las cláusulas de guardia (\textit{guard clauses}) minimizan la profundidad del árbol de bifurcaciones y optimizan la legibilidad del código.
\end{exampleblock}
\end{frame}

\begin{frame}{Caso de Estudio 2: Luminancia en Arte Digital}
    La percepción psicofísica humana del brillo se modela según la norma ITU-R BT.601:
    \begin{equation*}
        L = 0.299 \cdot R + 0.587 \cdot G + 0.114 \cdot B
    \end{equation*}
    \begin{block}{Segmentación Tonal Fotométrica}
        \begin{itemize}
            \item \textbf{Low-Key (Sombras):} $L < 64$
            \item \textbf{Mid-Tone (Tonos Medios):} $64 \le L < 192$
            \item \textbf{High-Key (Altas Luces):} $L \ge 192$
        \end{itemize}
    \end{block}
\end{frame}

\section{Evaluación Heurística: Matriz de Confusión}

\begin{frame}{La Matriz de Confusión Heurística}
    \begin{center}
    \begin{tabular}{c|cc}
        \toprule
        & \textbf{Predicho Positivo (Spam)} & \textbf{Predicho Negativo (Ham)} \\
        \midrule
        \textbf{Real Positivo} & \textbf{VP} (Ataque Mitigado) & \textbf{FN} (Brecha de Seguridad) \\
        \textbf{Real Negativo} & \textbf{FP} (Bloqueo Accidental) & \textbf{VN} (Tráfico Legítimo) \\
        \bottomrule
    \end{tabular}
    \end{center}
    \vspace{0.4cm}
    \begin{itemize}
        \item \textbf{Precisión:} $\frac{\text{VP}}{\text{VP} + \text{FP}}$ (Calidad de las alarmas emitidas).
        \item \textbf{Exhaustividad (\textit{Recall}):} $\frac{\text{VP}}{\text{VP} + \text{FN}}$ (Capacidad de interceptación total).
    \end{itemize}
\end{frame}

\begin{frame}{Conclusiones y Siguiente Paso}
    \begin{itemize}
        \item Las bifurcaciones \texttt{if-elif-else} constituyen la unidad atómica de decisión algorítmica.
        \item La evaluación de matrices de confusión permite auditar cuantitativamente reglas hechas a mano.
        \item \textbf{Siguiente Sesión (0.2):} Estructuras Multicamino con \texttt{match-case} y optimización temporal $\mathcal{O}(1)$ mediante Diccionarios de Despacho.
    \end{itemize}
\end{frame}

\end{document}
